\usepackage[T1]{fontenc}
\usepackage{lmodern}
\usepackage{geometry}
\usepackage{graphicx}
\usepackage{amsmath}
\usepackage{caption}
\usepackage{subcaption}
\usepackage{fancyhdr}
\usepackage{titlesec}
\usepackage{abstract}
%\usepackage{cite}
\usepackage{cancel}

\usepackage[dvipsnames]{xcolor}
%\usepackage[table]{xcolor}
\usepackage[most]{tcolorbox}
\tcbuselibrary{theorems}

\usepackage{amsthm}
\usepackage{amssymb}
\usepackage{pifont}
\usepackage{mathtools}
\usepackage{enumitem}
\usepackage{mathrsfs}
\usepackage{hyperref}
\usepackage{dirtytalk}
\usepackage{float}
\usepackage{comment}
\usepackage{nicematrix}
\usepackage{diagbox}
\usepackage{listings}

\NiceMatrixOptions{xdots/horizontal-labels}

\usepackage{algorithm}
\usepackage{algorithmic}
\usepackage{bussproofs}

	
%%% for figures
\usepackage{tikz}
\usetikzlibrary{decorations.markings}
\usetikzlibrary{graphs,graphdrawing,arrows.meta,decorations.pathreplacing}
\usegdlibrary{circular}
%\usegdlibrary{grid}  
\usetikzlibrary{positioning}
\usepackage{xcolor}
\usepackage[justification=centering]{caption}
\usepackage{xstring}
\usepackage{pgffor}
\usepackage{etoolbox}
%\setlength{\parskip}{2pt} 

%%% for bibliography
\usepackage{hyperref,url}
\renewcommand{\sectionautorefname}{Section}
\usepackage[numbers]{natbib}
\usepackage{xparse}

\NewDocumentCommand{\citehybridpage}{O{} m}{%
  [p.~#1]\cite{#2}%
}

\NewDocumentCommand{\citehybrid}{m}{%
  \cite{#1}
}


%\usepackage{cite}

\newcommand{\nrightsquigarrow}{
  \mathrel{\tikz[baseline=(a.base)]{
    \node[inner sep=0pt,outer sep=0pt] (a) {$\rightsquigarrow$};
    \draw  ([xshift=-0.3ex,yshift=-0.5ex]a.center) -- ([xshift=0.3ex,yshift=0.8ex]a.center);
  }}
}

\tcolorboxenvironment{theorem}{
  colback=black!7,       % light gray background
%%  colframe=black,        % frame color
%  boxrule=0pt,         % frame thickness
%  arc=0mm,               % rounded corners
%  left=6pt, right=6pt, top=6pt, bottom=6pt, % padding
%  enhanced,
  enhanced jigsaw,
  sharp corners,
%  colback=white,
  colframe=black!70,
  boxrule=0pt,
  borderline west={1pt}{0pt}{black!70},  % <-- vertical line on the left
  left=6pt,
  right=6pt,
  top=4pt,
  bottom=4pt,
  before skip=10pt,
  after skip=10pt,
}
\tcolorboxenvironment{notation}{
  colback=gray!10,       % light gray background
%  colframe=black,        % frame color
  boxrule=0pt,         % frame thickness
  arc=0mm,               % rounded corners
  left=6pt, right=6pt, top=6pt, bottom=6pt, % padding
  enhanced,
}
\tcolorboxenvironment{lemma}{
  colback=black!7,       % light gray background
%%  colframe=black,        % frame color
%  boxrule=0pt,         % frame thickness
%  arc=0mm,               % rounded corners
%  left=6pt, right=6pt, top=6pt, bottom=6pt, % padding
%  enhanced,
  enhanced jigsaw,
  sharp corners,
%  colback=white,
  colframe=black!70,
  boxrule=0pt,
  borderline west={1pt}{0pt}{black!70},  % <-- vertical line on the left
  left=6pt,
  right=6pt,
  top=4pt,
  bottom=4pt,
  before skip=10pt,
  after skip=10pt,
}
\tcolorboxenvironment{corollary}{
  colback=black!7,       % light gray background
%%  colframe=black,        % frame color
%  boxrule=0pt,         % frame thickness
%  arc=0mm,               % rounded corners
%  left=6pt, right=6pt, top=6pt, bottom=6pt, % padding
%  enhanced,
  enhanced jigsaw,
  sharp corners,
%  colback=white,
  colframe=black!70,
  boxrule=0pt,
  borderline west={1pt}{0pt}{black!70},  % <-- vertical line on the left
  left=6pt,
  right=6pt,
  top=4pt,
  bottom=4pt,
  before skip=10pt,
  after skip=10pt,
}
\tcolorboxenvironment{nameddefinition}{
  enhanced jigsaw,
  sharp corners,
  colback=white,
  colframe=black!70,
  boxrule=0pt,
  borderline west={1pt}{0pt}{black!70},  % <-- vertical line on the left
  left=6pt,
  right=6pt,
  top=4pt,
  bottom=4pt,
  before skip=10pt,
  after skip=10pt,
  fontupper=\slshape
}
 \tcolorboxenvironment{definition}{
  enhanced jigsaw,
  sharp corners,
  colback=white,
  colframe=black!70,
  boxrule=0pt,
  borderline west={1pt}{0pt}{black!70},  % <-- vertical line on the left
  left=6pt,
  right=6pt,
  top=4pt,
  bottom=4pt,
  before skip=10pt,
  after skip=10pt,
}
\tcolorboxenvironment{obs}{
  colback=Cyan!10,       % light gray background
%  colframe=black,        % frame color
  boxrule=0pt,         % frame thickness
  arc=0mm,               % rounded corners
  left=6pt, right=6pt, top=6pt, bottom=6pt, % padding
  enhanced,
}
\tcolorboxenvironment{proposition}{
  colback=black!7,       % light gray background
%%  colframe=black,        % frame color
%  boxrule=0pt,         % frame thickness
%  arc=0mm,               % rounded corners
%  left=6pt, right=6pt, top=6pt, bottom=6pt, % padding
%  enhanced,
  enhanced jigsaw,
  sharp corners,
%  colback=white,
  colframe=black!70,
  boxrule=0pt,
  borderline west={1pt}{0pt}{black!70},  % <-- vertical line on the left
  left=6pt,
  right=6pt,
  top=4pt,
  bottom=4pt,
  before skip=10pt,
  after skip=10pt,
}
\tcolorboxenvironment{remark}{
%  colback=WildStrawberry!10,       % light gray background
%%  colframe=black,        % frame color
%  boxrule=0pt,         % frame thickness
%  arc=0mm,               % rounded corners
%  left=6pt, right=6pt, top=6pt, bottom=6pt, % padding
%  enhanced,
  enhanced jigsaw,
  sharp corners,
  colback=white,
  colframe=black!70,
  boxrule=0pt,
  borderline west={1pt}{0pt}{black!70},  % <-- vertical line on the left
  left=6pt,
  right=6pt,
  top=4pt,
  bottom=4pt,
  before skip=10pt,
  after skip=10pt,
}

%\newtheoremstyle{slanted}
%  {\topsep}   % Space above
%  {\topsep}   % Space below
%  {\slshape}  % Body font (slanted)
%  {}          % Indent amount
%  {\bfseries} % Theorem head font
%  {.}         % Punctuation after theorem head
%  {.5em}      % Space after theorem head
%  {}          % Theorem head spec
%
%

\theoremstyle{definition}
\newtheorem{definition}{Definition}[section]
\newtheorem{nameddefinition}[definition]{Definition}

\newtheorem{theorem}{Theorem}[section]
\newtheorem{lemma}[theorem]{Lemma}
\newtheorem{notation}[theorem]{Notation}
\newtheorem{corollary}[theorem]{Corollary}
\newtheorem{obs}[theorem]{Observation}
\newtheorem{proposition}[theorem]{Proposition}
\newtheorem{remark}[theorem]{Remark}

\newtheorem{examplex}{Example}
\newenvironment{example}
  {\pushQED{\qed}\renewcommand{\qedsymbol}{$\lozenge$}\examplex}
  {\popQED\endexamplex}
	
\setlength{\columnsep}{20pt}

\geometry{margin=1.5cm}
%\usepackage[a4paper, left=2cm, right=2cm, bottom=2.54cm]{geometry}
\pagestyle{fancy}
\fancyhf{}
\rhead{\thepage}
\lhead{Automated consistency verification for replicated data systems}

\titleformat{\section}{\large\bfseries}{\thesection.}{1em}{}
\titleformat{\subsection}{\normalsize\bfseries}{\thesubsection.}{1em}{}



\usepackage[none]{hyphenat}
%%%%%%%%%%%%%%%%%%%%%%%%%%%%%%%%%%%%%%%%%%%%%%%%%%%%%%%%%%%%%%%%%%%%%%%%%%%%%%%%%%%%%%%%
\newcommand{\historyset}{\mathcal{H}(\mathbb{P}, \mathbb{T}, \mathbb{O}, \mathbb{V})}
\newcommand{\contexthistory}{Let $\mathbb{P}$, $\mathbb{T} = \{read,write\}$, $\mathbb{O}$, and $\mathbb{V}$ be finite sets representing, respectively, processes, operation types, objects, and values. }
\newcommand{\rb}{\xrightarrow{rb}}
\newcommand{\ssrel}{\xrightarrow{ss}}
\newcommand{\opa}{a}
\newcommand{\opb}{b}
\newcommand{\opc}{c}
\newcommand{\opd}{d}
\newcommand{\ope}{e}
\newcommand{\opf}{f}
\newcommand{\opg}{g}
\newcommand{\oph}{h}
\newcommand{\opi}{i}
\newcommand{\opj}{j}
\newcommand{\opk}{k}
\newcommand{\opl}{l}
\newcommand{\opm}{m}
\newcommand{\opn}{n}
\newcommand{\opo}{o}
\newcommand{\opp}{p}
\newcommand{\opq}{q}
\newcommand{\opr}{r}
\newcommand{\ops}{s}
\newcommand{\opt}{t}
\newcommand{\opu}{u}
\newcommand{\opv}{v}
\newcommand{\opw}{w}
\newcommand{\opx}{x}
\newcommand{\opy}{y}
\newcommand{\opz}{z}
\newcommand{\dsucp}[1]{\stackrel{#1}{\rightharpoonup}}
\newcommand{\dsuc}{\rightharpoonup}
%$\twoheadrightarrow$
\newcommand{\patharrow}{\!\mathrel{\rlap{\hskip 0.55em $\rightarrow$}\rightarrow}\;}
\newcommand{\algoneset}{\mathcal{G}_{\mathcal{A}\emph{lg1}}}
\newcommand{\algone}{\textsl{Alg1}}
\newcommand{\foralls}{\forall\;}
\newcommand{\ar}{\xrightarrow{ar}}
\newcommand{\vis}{\xrightarrow{vis}}

\makeatletter
\renewcommand{\@makefnmark}{\textsuperscript{[\thefootnote]}}
\makeatother
